\documentclass[11pt,a4paper]{article}

\usepackage[utf8]{inputenc}
\usepackage[T1]{fontenc}
\usepackage[margin=2.2cm]{geometry}
\usepackage{parskip}
\usepackage{titlesec}
\usepackage{enumitem}
\usepackage{booktabs}
\usepackage{xcolor}
\usepackage{tcolorbox}
\usepackage{hyperref}
\usepackage{fancyhdr}
\usepackage{listings}
\usepackage{longtable}
\usepackage{array}

\tcbuselibrary{skins,breakable}

\definecolor{accent}{RGB}{30,80,150}
\definecolor{softgray}{RGB}{245,246,248}
\definecolor{codebg}{RGB}{248,248,250}
\definecolor{warn}{RGB}{160,60,40}

\titleformat{\section}{\Large\bfseries\color{accent}}{\thesection}{0.8em}{}
\titleformat{\subsection}{\large\bfseries}{\thesubsection}{0.6em}{}
\titleformat{\subsubsection}{\normalsize\bfseries\itshape}{\thesubsubsection}{0.5em}{}

\hypersetup{
  colorlinks=true,
  linkcolor=accent,
  urlcolor=accent,
  pdftitle={AgentIA --- Current State},
  pdfauthor={AgentIA Engineering}
}

\pagestyle{fancy}
\fancyhf{}
\fancyhead[L]{\small\color{accent}AgentIA --- Current State}
\fancyhead[R]{\small\color{accent}\thepage}
\renewcommand{\headrulewidth}{0.4pt}

\lstset{
  basicstyle=\ttfamily\small,
  backgroundcolor=\color{codebg},
  frame=single,
  rulecolor=\color{softgray},
  breaklines=true,
  columns=fullflexible,
  keepspaces=true,
  showstringspaces=false,
  xleftmargin=0.5em,
  xrightmargin=0.5em,
}

\newtcolorbox{notebox}[1][]{colback=softgray,colframe=accent,boxrule=0.4pt,
  arc=2pt,left=6pt,right=6pt,top=4pt,bottom=4pt,breakable,#1}

\title{\vspace{-1.4cm}\textbf{AgentIA --- Current State}\\[0.25cm]
\large Bare bones: how it functions, what ships, and what is honestly not verified}
\author{AgentIA Engineering}
\date{\today}

\begin{document}
\maketitle
\thispagestyle{fancy}

\begin{notebox}
\textbf{Scope of this document.} The other three documents in \texttt{docs/} are
SkillOpt-focused and were written \emph{before} implementation:
\texttt{notes.tex} is a proposal brief, \texttt{resume\_skillopt.tex} is a method
reference with prompt templates, and \texttt{slides.tex} is a 13-frame deck. None
of them describes what is actually built. This document is that description, and it
is deliberately blunt about the gaps.
\end{notebox}

\tableofcontents
\vspace{0.6cm}

\section{What AgentIA is}

AgentIA --- \emph{Microservice Code Studio} --- ingests a formal specification of a
microservice and produces a compiled, tested, exportable Spring Boot project. It is
two layers that talk over REST and Server-Sent Events:

\begin{center}
\begin{tabular}{@{}p{0.44\textwidth}p{0.5\textwidth}@{}}
\toprule
\textbf{React / Vite Studio (port 3000)} & \textbf{FastAPI Orchestrator (port 8000)} \\
\midrule
Spec ingestion (BDD, architecture, relational models) &
FIFO concurrency queue (max 2 simultaneous builds) \\
Live generation monitor over SSE &
LangGraph autonomous agent (8 nodes) \\
Code explorer, security audit, DevOps, export &
Hermetic Docker sandbox (\texttt{mvn test -o --network none}) \\
ZIP download and atomic Git publish &
Bounded self-repair loop \\
& SQLite session store \\
\bottomrule
\end{tabular}
\end{center}

The property that makes it more than a code generator is that \textbf{it refuses to
claim success it did not verify}. That constraint was not present at the start; it
was built in specs 012 and 013, and it is the most important thing to understand
about the current state.

\section{How it functions}

\subsection{The LangGraph harness}

The agent is an eight-node graph. Topology as it exists in
\texttt{backend/app/orchestrator/graph.py}:

\begin{lstlisting}
START -> validator
validator  --(FAILED)--> END
           --(else)----> scaffolder
scaffolder -> domain -> service -> controller -> test -> sandbox
sandbox    --(build_success)--> END
           --(status BLOCKED)--> END
           --(else)-----------> repair
repair     --(status BLOCKED)--> END
           --(else)-----------> sandbox
\end{lstlisting}

\begin{center}
\begin{tabular}{@{}ll@{}}
\toprule
\textbf{Node} & \textbf{Responsibility} \\
\midrule
\texttt{validator} & Reject an unusable blueprint before any model call \\
\texttt{scaffolder} & \texttt{pom.xml}, \texttt{application.yml}, \texttt{@SpringBootApplication} \\
\texttt{domain} & JPA entities and DTO records \\
\texttt{service} & Service interfaces and implementations, repositories, exceptions \\
\texttt{controller} & REST controllers; error handling delegated to one advice class \\
\texttt{test} & One test per declared acceptance scenario \\
\texttt{sandbox} & Hermetic offline build and test, then the honesty check \\
\texttt{repair} & Bounded self-repair from parsed Maven output \\
\bottomrule
\end{tabular}
\end{center}

Two routing details matter and are easy to miss:

\begin{itemize}[leftmargin=1.4em]
  \item \texttt{sandbox -> END} on \texttt{BLOCKED}. An unverifiable workspace does
  \textbf{not} enter the repair loop: no code patch fixes a missing container
  runtime. This is what makes spec 012's honest failure terminate cleanly rather
  than burning the repair budget on an environment fault.
  \item \texttt{repair -> sandbox} is the loop, and it is bounded by the configured
  repair cap.
\end{itemize}

\subsection{The stage execution seam}

Every generation stage --- in both execution paths --- is invoked through one
function, \texttt{stages/runner.py::run\_stage}. It is the only place that owns mode
selection, request and correction budgeting, compliance gating, provenance
recording, cost recording, and journal accumulation. That the two paths share it is
what makes behavioural parity structural rather than a matter of discipline.

\begin{center}
\begin{tabular}{@{}p{0.16\textwidth}p{0.78\textwidth}@{}}
\toprule
\textbf{Path} & \textbf{Entry} \\
\midrule
A --- graph & \texttt{generation\_graph.stream} via
\texttt{routes\_session.execute\_generation\_pipeline}; nodes are thin delegations
into the seam \\
B --- sequential & \texttt{pipeline\_runner.\_execute\_pipeline\_steps} (auto-pilot),
calling the same seam directly \\
\bottomrule
\end{tabular}
\end{center}

\subsection{The two generation modes}

A session is decided to be \texttt{MODEL} or \texttt{DETERMINISTIC} once, at start,
and never changes. \texttt{MODEL} never falls back to deterministic.

\begin{center}
\begin{tabular}{@{}lll@{}}
\toprule
& \textbf{DETERMINISTIC} & \textbf{MODEL} \\
\midrule
Implementation & \texttt{stages/deterministic/*.emit} & \texttt{stages/model/*} \\
Compliance gate & none & applied before persistence \\
Model requests & 0 & 1 initial $+$ up to 2 corrections per stage \\
Provenance & no provider or model recorded & provider, model, revision, ordinal \\
\bottomrule
\end{tabular}
\end{center}

The gate asymmetry is deliberate (FR-013): deterministic output is compliant by
construction, and gating it would change offline behaviour. It is documented at the
point of use so it is not later ``fixed'' as an oversight.

\subsection{The compliance gate}

Two independently-written validator families evaluate every candidate artifact set.
They overlap and \textbf{disagree on severity} for the same violation --- Lombok
\texttt{@Data} on a JPA entity is MEDIUM in one and HIGH in the other. The adapter
keeps the strictest, and that tightening is a real behaviour change: it can convert
previously-completing sessions into blocked ones. It is recorded in
\texttt{specs/011-llm-generation-nodes/migration-notes.md} because it is the single
most likely source of a false regression attribution.

Violations are attributed one of two ways, and the distinction is load-bearing:

\begin{center}
\begin{tabular}{@{}lll@{}}
\toprule
\textbf{Attribution} & \textbf{Effect} & \textbf{Canonical case} \\
\midrule
LOCAL & Rejects the candidate set; consumes a correction attempt &
DTO declared as a class \\
ACCUMULATED & Recorded, never rejects, never consumes an attempt &
No \texttt{@RestControllerAdvice} anywhere in the project \\
\bottomrule
\end{tabular}
\end{center}

The accumulated case is why a scaffolder running \emph{before} the controller stage
is not rejected for a missing global error handler it could not have written.

\subsection{The honest verifier (spec 012)}

The sandbox originally substituted a synthetic \texttt{BUILD SUCCESS} with
fabricated 5/5 test counts whenever it could not build. Sessions therefore reached
\texttt{VERIFIED} for workspaces that were never compiled, and every downstream
figure inherited that. It is now:

\begin{center}
\begin{tabular}{@{}lll@{}}
\toprule
\textbf{Outcome} & \texttt{exit\_code} & \texttt{fallback\_used} \\
\midrule
Verified and passed & 0 & false \\
Verified and failed & non-zero & false \\
\textbf{Could not verify} & non-zero & \textbf{true} \\
\bottomrule
\end{tabular}
\end{center}

There are \textbf{four} substitution triggers, not the three originally identified:
runtime unavailable, runtime executable missing, runtime communication failure, and
--- the most severe --- a build failure whose output matches an environment-looking
pattern, which used to be converted into a \emph{pass}.

An unverifiable session terminates in \texttt{BLOCKED}, bypassing repair, with the
reason recorded. \texttt{ALLOW\_HERMETIC\_FALLBACK} (default \textbf{false})
restores the old behaviour for local development; the marking is recorded in
\textbf{both} modes, because permissive mode changes what is permitted, not what is
recorded.

\subsection{Cost tracing (spec 013)}

Every model call is recorded through a proxy applied by the model factory --- but
only while a recording context is active, which is true inside the stage boundary
and false for a direct factory call. That condition keys on \emph{is recording
active}, never on \emph{is this a test}; a test-shaped bypass would leave production
unrecorded while the suite stayed green.

The local SQLite store is the \textbf{system of record}. The telemetry destination is
a mirror whose absence is a non-event, so the report is deterministic and offline.
Pricing is per call, not per model, because the provider's published rates vary by
peak/off-peak (2$\times$) and by cache hit/miss (up to 100$\times$ on input).

\subsection{The SkillOpt loop (spec 014)}

One iteration, one skill, no memory between runs:

\begin{lstlisting}
collect recorded outcomes -> reflect into <=4 edits -> apply to a COPY
  -> gate BOTH skills on the SAME fresh executions
  -> accept iff candidate_score > current_score   (strict; ties rejected)
  -> log exactly one run record
\end{lstlisting}

Every stability mechanism the source method uses is deliberately absent: no
rejected-edit buffer, no learning-rate schedule, no epoch-wise slow/meta update, no
multi-skill bank, no significance test. The protected
\texttt{SLOW\_UPDATE} region is already honoured so that adding the slow update later
needs no change to the edit contract.

\section{Features shipped}

\begin{longtable}{@{}p{0.06\textwidth}p{0.30\textwidth}p{0.56\textwidth}@{}}
\toprule
\textbf{Spec} & \textbf{Feature} & \textbf{State} \\
\midrule
\endhead
001 & Microservice Code Studio & Core platform: ingestion, queue, session store, SSE \\
002 & Requirements to stories & BDD ingestion to Given/When/Then \\
003 & Architecture component design & Component and Mermaid rendering \\
004 & Domain models / SQL & Relational model generation \\
005 & Automated tests analysis & Test suite analysis and repair histories \\
006 & Security \& quality audit & SAST, quality gate, export blocking \\
007 & Docker CI/CD orchestration & Deployment and pipeline surfaces \\
008 & Unified workflow orchestration & Lifecycle, phases, auto-pilot \\
009 & Historical baseline analysis & Synthetic incident baseline \\
010 & Skill injection & \textbf{Domain triage only} --- the skill documents were never authored \\
011 & LLM generation nodes & \textbf{Built.} Seam, 5 model stages, gate, correction loop, provenance, journal, measurement harness \\
012 & Sandbox verifier honesty & \textbf{Built.} Four triggers, permissive opt-in, terminal-state parity \\
013 & Cost tracing & \textbf{Built.} Per-call telemetry, aggregation, report, pricing table \\
014 & SkillOpt compact & \textbf{Built.} One-iteration loop, strict gate, disjoint evidence \\
\bottomrule
\end{longtable}

\begin{notebox}
\textbf{Spec 010 is the notable gap.} It classified the failure domains and created
three skill files, but all three are \textbf{zero bytes} --- placeholders. Only
\texttt{layer\_architecture.md} has content today, authored as feature 014's seed.
\texttt{exception\_handling.md} and \texttt{mockito\_tests.md} remain empty.
\end{notebox}

\section{The agent scripts}

Everything under \texttt{backend/scripts/} is run directly; none is wired into the
API.

\subsection{Generation and measurement}

\begin{longtable}{@{}p{0.36\textwidth}p{0.60\textwidth}@{}}
\toprule
\textbf{Script} & \textbf{Purpose} \\
\midrule
\endhead
\texttt{capture\_generation\_baseline.py} & Freezes the pre-migration deterministic
output (spec 011 T008). Must run \emph{before} any stage is migrated. \\
\texttt{analyze\_historical\_baseline.py} & Analyses the synthetic incident baseline
(spec 009) \\
\texttt{seed\_historical\_baseline.py} & Seeds it \\
\texttt{report\_session\_costs.py} & \textbf{The cost report.} Prints the headline
average cost per completed session with its population, exclusions and pricing basis \\
\texttt{validate\_seed\_skills.py} & Validates skill documents \\
\bottomrule
\end{longtable}

\subsection{SkillOpt}

\begin{longtable}{@{}p{0.30\textwidth}p{0.66\textwidth}@{}}
\toprule
\textbf{Module} & \textbf{Role} \\
\midrule
\endhead
\texttt{run\_skillopt.py} & The orchestrator: one iteration, one run record \\
\texttt{skillopt/collect.py} & Recorded sessions to outcome records. ``No failures''
$\neq$ ``no sessions''; a fallback-marked session counts as a failure \\
\texttt{skillopt/reflect.py} & One direct factory call to a bounded edit set \\
\texttt{skillopt/apply.py} & Four operations, cap $L_t = 4$, applied to a
\textbf{copy}; protected-region and not-found targets rejected individually \\
\texttt{skillopt/gate.py} & Both skills on the same fresh executions; strict
acceptance; disjointness asserted in code \\
\texttt{skillopt/real\_runner.py} & The production session runner, used only by
\texttt{--real} \\
\texttt{skillopt/prompts/analyst\_error.md} & The reflector prompt, adapted from
SkillOpt Appendix C.2.1 \\
\bottomrule
\end{longtable}

\section{Economics: where the numbers live}

\begin{longtable}{@{}p{0.42\textwidth}p{0.54\textwidth}@{}}
\toprule
\textbf{What} & \textbf{Where} \\
\midrule
\endhead
The headline report & \texttt{backend/scripts/report\_session\_costs.py} \\
System of record & \texttt{backend/cost\_tracking.db} --- tables
\texttt{model\_calls} and \texttt{session\_costs} \\
Published rate table & \texttt{backend/app/resources/model\_pricing.json} (cites its
source URL in-file) \\
Intervention rate & \texttt{reports/measurements/011-sc011-intervention-rate.md} \\
Excluded-session count & \texttt{reports/measurements/012-fallback-excluded-sessions.md} \\
Frozen pre-migration baseline & \texttt{reports/baselines/} --- immutable \\
\bottomrule
\end{longtable}

Producing the number:

\begin{lstlisting}
PYTHONPATH=backend .venv/bin/python backend/scripts/report_session_costs.py
\end{lstlisting}

\begin{notebox}[colframe=warn]
\textbf{There is no cost figure yet, and reporting one would be the failure this
platform was built to prevent.} The cost store is currently \textbf{empty}: no
session has run with recording active. The report says so explicitly ---
\emph{``No sessions found\ldots This is not a measurement of zero cost''} --- and
that refusal is the correct behaviour, not a missing feature.

The only quantified figure that exists is the SC-011 intervention rate:
\textbf{12.5\% over 32 sessions}, and it is labelled \textbf{synthetic} in the
artifact itself, because those sessions were driven by a fake client with a
scripted compliant/non-compliant mix. It measures the mechanism, not a real model.
\end{notebox}

\section{How to test it}

\subsection{Start both halves}

The \texttt{.bat} and \texttt{.ps1} launchers assume Windows and a bare
\texttt{python}. On Linux:

\begin{lstlisting}
# backend  -> http://localhost:8000   (Swagger UI at /docs)
cd backend && ../.venv/bin/python -m uvicorn app.main:app --reload --port 8000

# frontend -> http://localhost:3000
cd frontend && npm run dev
\end{lstlisting}

\texttt{frontend/src/services/apiClient.ts} uses \texttt{baseURL: '/api/v1'}, so
the Vite dev server needs a proxy to port 8000.

\subsection{One run that exercises the newest machinery}

A single \texttt{MODEL}-mode session on the \texttt{minimal} blueprint touches the
stage seam (011), the honest verifier (012) and cost recording (013).

\begin{center}
\begin{tabular}{@{}clp{0.52\textwidth}@{}}
\toprule
\textbf{\#} & \textbf{View} & \textbf{What to observe} \\
\midrule
1 & \texttt{SpecIngestionView} & Paste
\texttt{backend/tests/fixtures/baseline\_blueprints/minimal.json}; provider
\texttt{deepseek} $+$ key \\
2 & \texttt{GenerationMonitorView} & Live SSE: phase transitions, build logs, and
the \texttt{verification\_result} event carrying \texttt{verificationFallbackUsed} \\
3 & \texttt{CodeExplorerView} & Generated artifacts and their provenance \\
4 & \texttt{ExportPublishView} & The ZIP. A \textbf{403} means the quality gate
blocked the export --- correct, not a bug \\
\bottomrule
\end{tabular}
\end{center}

API equivalent (request shape verified against
\texttt{QuickStartSessionRequest}):

\begin{lstlisting}
curl -X POST http://localhost:8000/api/v1/sessions/quick-start \
  -H 'Content-Type: application/json' \
  -d '{"serviceName":"notes-service","rawText":"<minimal.json>",
       "autoRun":true,"llmProvider":"deepseek","apiKey":"sk-..."}'

# then:
#   GET /api/v1/sessions/{id}/stream      SSE
#   GET /api/v1/sessions/{id}             detail (+ verificationFallbackUsed)
#   GET /api/v1/sessions/{id}/artifacts   listing
\end{lstlisting}

\subsection{Two host-dependent behaviours to expect}

\textbf{The sandbox result depends on container-runtime reachability.} If the
runtime is unreachable and \texttt{ALLOW\_HERMETIC\_FALLBACK} is unset, the session
\textbf{blocks}. That is spec 012 working. To run the loop end to end on such a host,
either enable the permissive flag deliberately or make the runtime reachable.

\textbf{The SkillOpt real iteration needs two flags, not one.} A key in
\texttt{backend/.env} reaches \texttt{os.environ} through \texttt{load\_dotenv()},
so a key-presence check passes without the operator realising a key exists. It
therefore requires \texttt{AGENTIA\_RUN\_REAL\_SKILLOPT=1} \emph{and}
\texttt{AGENTIA\_ALLOW\_REAL\_SPEND=1}, and it spends real money.

\section{Honest current limits}

Recorded here rather than left to be discovered.

\begin{center}
\begin{tabular}{@{}p{0.30\textwidth}p{0.64\textwidth}@{}}
\toprule
\textbf{Limit} & \textbf{Detail} \\
\midrule
No cost figure exists & The store is empty. Running the report is the correct first
step; reporting a number before that is the failure mode 012/013 close. \\
SC-003 not verified by the implementer & The real-build criterion was verified by
the operator from a capable shell; the implementer's shell cannot reach the
container runtime. Recorded as such, never inferred. \\
SC-007 partially observed & A real SkillOpt iteration ran twice unintentionally and
reached a decision, but its scores were suppressed by the test and were not
re-run, to avoid spending again without authorisation. \\
The gate is coarse & Five blueprints, $M = 4$: one task moves the pass rate by
\textbf{25 points}. Noise can decide an acceptance. \\
No memory between runs & A rejected edit is logged and dropped, so v1 can repeat a
failed edit indefinitely and nothing notices. \\
Repair-cap drift & The constitution specifies 3; \texttt{config.py} sets
\textbf{5}. This flatters any absolute intervention-rate ceiling and is disclosed in
the measurement artifact. \\
Three skill files empty & \texttt{exception\_handling.md} and
\texttt{mockito\_tests.md} are zero bytes; only the seed has content. \\
No frontend surface for cost & The report is a CLI artifact. There is no view for
it. \\
\bottomrule
\end{tabular}
\end{center}

\section{Next steps}

\subsection{Before anything else}

\textbf{Widen the held-out exam.} Adding optimizer machinery to a five-task gate
optimises against noise. This is the highest-value change and it is cheap.

\textbf{Run sessions with recording on.} The economics are unmeasurable until the
store has data, and every other cost question is downstream of that.

\subsection{SkillOpt design improvements}

\begin{enumerate}[leftmargin=1.4em]
  \item \textbf{Filter environment failures in code, not in the prompt.} The
  collector currently classifies fallback-marked sessions as failures, sends them to
  the reflector, and the prompt then asks the model to ignore them. Keeping them for
  context is defensible, but the ignore-instruction should not be the only guard.
  \item \textbf{Cache the current score} keyed by (skill hash, held-out set). The
  gate is $2 \times M$ real sessions per iteration; caching makes an epoch loop
  affordable at all.
  \item \textbf{Record a real build exit code.} It is currently \emph{derived} from
  the verification metrics' pass flag because the session table has no such column.
  \item \textbf{Pass provider and key into the session state} in
  \texttt{real\_runner}; a real run currently depends on ambient environment
  detection.
  \item \textbf{Then, in the source method's order:} rejected-edit buffer $\to$
  textual learning-rate schedule $\to$ epoch/slow-meta update $\to$ multi-skill
  bank. The protected region already exists for the slow update.
  \item \textbf{Statistical validation} --- the deferred item that makes any of the
  above trustworthy at this sample size.
\end{enumerate}

\subsection{Documentation}

\texttt{README.md} predates specs 011--014 entirely: it still describes the agent as
blocks that no longer exist, and states a repair cap of 3 when the code sets 5. It
should either point here or be brought current.

\end{document}
